% Options for packages loaded elsewhere
\PassOptionsToPackage{unicode}{hyperref}
\PassOptionsToPackage{hyphens}{url}
%
\documentclass[
]{article}
\usepackage{amsmath,amssymb}
\usepackage{iftex}
\ifPDFTeX
  \usepackage[T1]{fontenc}
  \usepackage[utf8]{inputenc}
  \usepackage{textcomp} % provide euro and other symbols
\else % if luatex or xetex
  \usepackage{unicode-math} % this also loads fontspec
  \defaultfontfeatures{Scale=MatchLowercase}
  \defaultfontfeatures[\rmfamily]{Ligatures=TeX,Scale=1}
\fi
\usepackage{lmodern}
\ifPDFTeX\else
  % xetex/luatex font selection
\fi
% Use upquote if available, for straight quotes in verbatim environments
\IfFileExists{upquote.sty}{\usepackage{upquote}}{}
\IfFileExists{microtype.sty}{% use microtype if available
  \usepackage[]{microtype}
  \UseMicrotypeSet[protrusion]{basicmath} % disable protrusion for tt fonts
}{}
\makeatletter
\@ifundefined{KOMAClassName}{% if non-KOMA class
  \IfFileExists{parskip.sty}{%
    \usepackage{parskip}
  }{% else
    \setlength{\parindent}{0pt}
    \setlength{\parskip}{6pt plus 2pt minus 1pt}}
}{% if KOMA class
  \KOMAoptions{parskip=half}}
\makeatother
\usepackage{xcolor}
\usepackage[margin=1in]{geometry}
\usepackage{color}
\usepackage{fancyvrb}
\newcommand{\VerbBar}{|}
\newcommand{\VERB}{\Verb[commandchars=\\\{\}]}
\DefineVerbatimEnvironment{Highlighting}{Verbatim}{commandchars=\\\{\}}
% Add ',fontsize=\small' for more characters per line
\usepackage{framed}
\definecolor{shadecolor}{RGB}{248,248,248}
\newenvironment{Shaded}{\begin{snugshade}}{\end{snugshade}}
\newcommand{\AlertTok}[1]{\textcolor[rgb]{0.94,0.16,0.16}{#1}}
\newcommand{\AnnotationTok}[1]{\textcolor[rgb]{0.56,0.35,0.01}{\textbf{\textit{#1}}}}
\newcommand{\AttributeTok}[1]{\textcolor[rgb]{0.13,0.29,0.53}{#1}}
\newcommand{\BaseNTok}[1]{\textcolor[rgb]{0.00,0.00,0.81}{#1}}
\newcommand{\BuiltInTok}[1]{#1}
\newcommand{\CharTok}[1]{\textcolor[rgb]{0.31,0.60,0.02}{#1}}
\newcommand{\CommentTok}[1]{\textcolor[rgb]{0.56,0.35,0.01}{\textit{#1}}}
\newcommand{\CommentVarTok}[1]{\textcolor[rgb]{0.56,0.35,0.01}{\textbf{\textit{#1}}}}
\newcommand{\ConstantTok}[1]{\textcolor[rgb]{0.56,0.35,0.01}{#1}}
\newcommand{\ControlFlowTok}[1]{\textcolor[rgb]{0.13,0.29,0.53}{\textbf{#1}}}
\newcommand{\DataTypeTok}[1]{\textcolor[rgb]{0.13,0.29,0.53}{#1}}
\newcommand{\DecValTok}[1]{\textcolor[rgb]{0.00,0.00,0.81}{#1}}
\newcommand{\DocumentationTok}[1]{\textcolor[rgb]{0.56,0.35,0.01}{\textbf{\textit{#1}}}}
\newcommand{\ErrorTok}[1]{\textcolor[rgb]{0.64,0.00,0.00}{\textbf{#1}}}
\newcommand{\ExtensionTok}[1]{#1}
\newcommand{\FloatTok}[1]{\textcolor[rgb]{0.00,0.00,0.81}{#1}}
\newcommand{\FunctionTok}[1]{\textcolor[rgb]{0.13,0.29,0.53}{\textbf{#1}}}
\newcommand{\ImportTok}[1]{#1}
\newcommand{\InformationTok}[1]{\textcolor[rgb]{0.56,0.35,0.01}{\textbf{\textit{#1}}}}
\newcommand{\KeywordTok}[1]{\textcolor[rgb]{0.13,0.29,0.53}{\textbf{#1}}}
\newcommand{\NormalTok}[1]{#1}
\newcommand{\OperatorTok}[1]{\textcolor[rgb]{0.81,0.36,0.00}{\textbf{#1}}}
\newcommand{\OtherTok}[1]{\textcolor[rgb]{0.56,0.35,0.01}{#1}}
\newcommand{\PreprocessorTok}[1]{\textcolor[rgb]{0.56,0.35,0.01}{\textit{#1}}}
\newcommand{\RegionMarkerTok}[1]{#1}
\newcommand{\SpecialCharTok}[1]{\textcolor[rgb]{0.81,0.36,0.00}{\textbf{#1}}}
\newcommand{\SpecialStringTok}[1]{\textcolor[rgb]{0.31,0.60,0.02}{#1}}
\newcommand{\StringTok}[1]{\textcolor[rgb]{0.31,0.60,0.02}{#1}}
\newcommand{\VariableTok}[1]{\textcolor[rgb]{0.00,0.00,0.00}{#1}}
\newcommand{\VerbatimStringTok}[1]{\textcolor[rgb]{0.31,0.60,0.02}{#1}}
\newcommand{\WarningTok}[1]{\textcolor[rgb]{0.56,0.35,0.01}{\textbf{\textit{#1}}}}
\usepackage{graphicx}
\makeatletter
\def\maxwidth{\ifdim\Gin@nat@width>\linewidth\linewidth\else\Gin@nat@width\fi}
\def\maxheight{\ifdim\Gin@nat@height>\textheight\textheight\else\Gin@nat@height\fi}
\makeatother
% Scale images if necessary, so that they will not overflow the page
% margins by default, and it is still possible to overwrite the defaults
% using explicit options in \includegraphics[width, height, ...]{}
\setkeys{Gin}{width=\maxwidth,height=\maxheight,keepaspectratio}
% Set default figure placement to htbp
\makeatletter
\def\fps@figure{htbp}
\makeatother
\setlength{\emergencystretch}{3em} % prevent overfull lines
\providecommand{\tightlist}{%
  \setlength{\itemsep}{0pt}\setlength{\parskip}{0pt}}
\setcounter{secnumdepth}{-\maxdimen} % remove section numbering
\ifLuaTeX
  \usepackage{selnolig}  % disable illegal ligatures
\fi
\IfFileExists{bookmark.sty}{\usepackage{bookmark}}{\usepackage{hyperref}}
\IfFileExists{xurl.sty}{\usepackage{xurl}}{} % add URL line breaks if available
\urlstyle{same}
\hypersetup{
  pdftitle={Share of Search as a complement to brand tracking},
  pdfauthor={Zen},
  hidelinks,
  pdfcreator={LaTeX via pandoc}}

\title{Share of Search as a complement to brand tracking}
\usepackage{etoolbox}
\makeatletter
\providecommand{\subtitle}[1]{% add subtitle to \maketitle
  \apptocmd{\@title}{\par {\large #1 \par}}{}{}
}
\makeatother
\subtitle{A pilot study on 50 Australian categories --- Tracksuit DS
take-home}
\author{Zen}
\date{16 August 2026}

\begin{document}
\maketitle

{
\setcounter{tocdepth}{2}
\tableofcontents
}
\hypertarget{summary-for-non-technical-readers}{%
\section*{Summary for non-technical
readers}\label{summary-for-non-technical-readers}}
\addcontentsline{toc}{section}{Summary for non-technical readers}

\textbf{What we did.} We collected 3.5 years of Google search interest
for every brand Tracksuit tracks in 50 Australian categories and
compared each brand's \emph{Share of Search} (its slice of all searches
for brands in its category) against the survey funnel metrics ---
awareness, consideration, preference.

\textbf{What we found.}

Search ranks brands the way the survey does --- in the right categories.
Where brands have distinctive names and consumers genuinely search for
them, Share of Search reproduces the survey's brand ranking closely.
Where they don't (small FMCG brands, generic names), roughly a third of
brands are simply invisible in search data.

Search cannot yet replace the survey for tracking \emph{movement}.
Month-to-month changes in consideration are mostly survey sampling noise
at n≈200 --- and we show that this, not search quality, is what caps any
correlation. Search may in fact be the \emph{steadier} signal in
low-sample categories.

Search shows a hint of \emph{leading} consideration by two to three
months (suggestive under a placebo test, not yet established) ---
consistent with published evidence that Share of Search leads market
share by 6--12 months. If it holds up under further validation, this is
the most commercially interesting property: an early-warning indicator
between survey waves.

\textbf{What we recommend.} Use search as a \emph{complement}, not a
substitute: launch a Share-of-Search module in categories that pass a
measurable eligibility test, position it for competitive benchmarking
and early trend detection, and do not position it as a replacement for
survey-based funnel tracking --- search carries no demographics, no
funnel stages, and no absolute levels.

\emph{(Caveats: Google Trends is a relative, sampled index;
brand-to-search-term mapping requires curation --- we quantify how much
it matters; Australia only; one search source.)}

\hypertarget{the-question}{%
\section{1. The question}\label{the-question}}

Share of Search (SoS) is increasingly cited as a cheap proxy for brand
strength. But ``does search correlate with our metrics?'' is really two
different product questions:

\begin{enumerate}
\def\labelenumi{\arabic{enumi}.}
\tightlist
\item
  \textbf{Between brands:} does SoS rank and size brands within a
  category the way our survey does? If yes, search could
  \emph{substitute} for parts of competitive benchmarking.
\item
  \textbf{Within a brand over time:} do movements in SoS track (or lead)
  movements in that brand's consideration? If yes, search could
  \emph{complement} the survey as a higher-frequency trend signal.
\end{enumerate}

These have different bars to clear and different failure modes, so we
test them separately, on the same data.

Search is treated throughout as a proxy for \textbf{consideration} (both
imply intent --- one behavioral, one claimed), but we also test which
funnel metric search actually resembles most (§6).

\hypertarget{data}{%
\section{2. Data}\label{data}}

\hypertarget{survey-data}{%
\subsection{2.1 Survey data}\label{survey-data}}

\begin{Shaded}
\begin{Highlighting}[]
\NormalTok{survey }\SpecialCharTok{|\textgreater{}}
  \FunctionTok{summarise}\NormalTok{(}\AttributeTok{categories =} \FunctionTok{n\_distinct}\NormalTok{(category\_name),}
            \StringTok{\textasciigrave{}}\AttributeTok{brand{-}category pairs}\StringTok{\textasciigrave{}} \OtherTok{=} \FunctionTok{n\_distinct}\NormalTok{(}\FunctionTok{paste}\NormalTok{(category\_name, brand\_name)),}
            \AttributeTok{waves =} \FunctionTok{n\_distinct}\NormalTok{(wave\_date),}
            \AttributeTok{first\_wave =} \FunctionTok{min}\NormalTok{(wave\_date), }\AttributeTok{last\_wave =} \FunctionTok{max}\NormalTok{(wave\_date))}
\end{Highlighting}
\end{Shaded}

\begin{verbatim}
## # A tibble: 1 x 5
##   categories `brand-category pairs` waves first_wave last_wave 
##        <int>                  <int> <int> <date>     <date>    
## 1         50                    507    43 2021-09-01 2025-03-01
\end{verbatim}

Monthly syndicated tracking, \textasciitilde200 qualified respondents
per category-wave. Funnel metrics: prompted awareness, consideration,
preference.

\hypertarget{search-data}{%
\subsection{2.2 Search data}\label{search-data}}

Google Trends, Australia, monthly, September 2021 -- March 2025,
collected by \texttt{pull/pull\_trends.py} (see
\texttt{pull/README.md}):

\begin{itemize}
\tightlist
\item
  Brands are resolved to Google \textbf{knowledge-graph topic entities}
  where a credible one exists (the main defense against homonyms ---
  Koala the mattress vs the animal, Mecca the retailer vs the city),
  with hand-curated terms for the pilot categories and every resolution
  logged for review.
\item
  Trends only compares 5 terms at a time on a relative 0--100 scale, so
  each category's highest-awareness brand rides along in every request
  as an \textbf{anchor}, and batches are rescaled to the anchor's level
  --- putting all brands in a category on one common scale.
\item
  \textbf{Share of Search} = a brand's index / sum of all its category's
  brand indices, per month.
\end{itemize}

\begin{Shaded}
\begin{Highlighting}[]
\NormalTok{trends }\SpecialCharTok{|\textgreater{}}
  \FunctionTok{summarise}\NormalTok{(}\AttributeTok{categories =} \FunctionTok{n\_distinct}\NormalTok{(category\_name),}
            \StringTok{\textasciigrave{}}\AttributeTok{brand{-}category pairs}\StringTok{\textasciigrave{}} \OtherTok{=} \FunctionTok{n\_distinct}\NormalTok{(}\FunctionTok{paste}\NormalTok{(category\_name, brand\_name)),}
            \AttributeTok{months =} \FunctionTok{n\_distinct}\NormalTok{(month))}
\end{Highlighting}
\end{Shaded}

\begin{verbatim}
## # A tibble: 1 x 3
##   categories `brand-category pairs` months
##        <int>                  <int>  <int>
## 1         50                    507     44
\end{verbatim}

\begin{Shaded}
\begin{Highlighting}[]
\NormalTok{vol }\OtherTok{\textless{}{-}}\NormalTok{ trends }\SpecialCharTok{|\textgreater{}}
  \FunctionTok{group\_by}\NormalTok{(category\_name, brand\_name) }\SpecialCharTok{|\textgreater{}}
  \FunctionTok{summarise}\NormalTok{(}\AttributeTok{mean\_idx =} \FunctionTok{mean}\NormalTok{(trends\_index, }\AttributeTok{na.rm =} \ConstantTok{TRUE}\NormalTok{), }\AttributeTok{.groups =} \StringTok{"drop"}\NormalTok{)}
\NormalTok{searchable\_share }\OtherTok{\textless{}{-}} \FunctionTok{mean}\NormalTok{(vol}\SpecialCharTok{$}\NormalTok{mean\_idx }\SpecialCharTok{\textgreater{}=}\NormalTok{ VOL\_MIN)}
\end{Highlighting}
\end{Shaded}

\textbf{A finding before any correlation:} only 81\% of tracked brands
register measurable search volume at all. The long tail of
survey-tracked brands is invisible to Google Trends at national scale.
This alone rules out search as a full survey substitute, and it becomes
the first criterion of the eligibility framework (§8). Brands below the
volume floor are excluded from SoS computations below.

\hypertarget{noise}{%
\section{3. How much signal is even in the survey?}\label{noise}}

Before correlating anything we ask what a correlation against monthly
survey metrics \emph{could possibly} achieve. A monthly consideration
reading rests on \textasciitilde200 qualified respondents, so it carries
sampling error of roughly ±2--3.5pp. Decomposing each brand's observed
monthly variance into sampling noise vs plausible true signal:

\begin{Shaded}
\begin{Highlighting}[]
\NormalTok{nd }\OtherTok{\textless{}{-}} \FunctionTok{noise\_decomposition}\NormalTok{(survey, }\StringTok{"consideration"}\NormalTok{)}
\NormalTok{nd\_cat }\OtherTok{\textless{}{-}}\NormalTok{ nd }\SpecialCharTok{|\textgreater{}} \FunctionTok{filter}\NormalTok{(waves }\SpecialCharTok{\textgreater{}=} \DecValTok{12}\NormalTok{) }\SpecialCharTok{|\textgreater{}}
  \FunctionTok{group\_by}\NormalTok{(category\_name) }\SpecialCharTok{|\textgreater{}}
  \FunctionTok{summarise}\NormalTok{(}\AttributeTok{brands =} \FunctionTok{n}\NormalTok{(), }\AttributeTok{mean\_reliability =} \FunctionTok{mean}\NormalTok{(reliability, }\AttributeTok{na.rm =} \ConstantTok{TRUE}\NormalTok{),}
            \AttributeTok{mean\_max\_r =} \FunctionTok{mean}\NormalTok{(max\_attainable\_r, }\AttributeTok{na.rm =} \ConstantTok{TRUE}\NormalTok{), }\AttributeTok{.groups =} \StringTok{"drop"}\NormalTok{)}

\NormalTok{median\_se }\OtherTok{\textless{}{-}}\NormalTok{ survey }\SpecialCharTok{|\textgreater{}} \FunctionTok{mutate}\NormalTok{(}\AttributeTok{se =} \FunctionTok{se\_prop}\NormalTok{(consideration, base\_n)) }\SpecialCharTok{|\textgreater{}}
  \FunctionTok{summarise}\NormalTok{(}\AttributeTok{m =} \FunctionTok{median}\NormalTok{(se, }\AttributeTok{na.rm =} \ConstantTok{TRUE}\NormalTok{)) }\SpecialCharTok{|\textgreater{}} \FunctionTok{pull}\NormalTok{(m)}

\FunctionTok{ggplot}\NormalTok{(nd\_cat, }\FunctionTok{aes}\NormalTok{(mean\_reliability, }\FunctionTok{reorder}\NormalTok{(category\_name, mean\_reliability))) }\SpecialCharTok{+}
  \FunctionTok{geom\_col}\NormalTok{(}\AttributeTok{fill =}\NormalTok{ PAL[}\DecValTok{1}\NormalTok{], }\AttributeTok{width =}\NormalTok{ .}\DecValTok{75}\NormalTok{) }\SpecialCharTok{+}
  \FunctionTok{scale\_x\_continuous}\NormalTok{(}\AttributeTok{labels =}\NormalTok{ scales}\SpecialCharTok{::}\NormalTok{percent, }\AttributeTok{expand =} \FunctionTok{expansion}\NormalTok{(}\AttributeTok{mult =} \FunctionTok{c}\NormalTok{(}\DecValTok{0}\NormalTok{, .}\DecValTok{05}\NormalTok{))) }\SpecialCharTok{+}
  \FunctionTok{labs}\NormalTok{(}\AttributeTok{x =} \StringTok{"Share of monthly variance in consideration that is true signal"}\NormalTok{,}
       \AttributeTok{y =} \ConstantTok{NULL}\NormalTok{,}
       \AttributeTok{title =} \StringTok{"Most monthly movement in consideration is sampling noise"}\NormalTok{,}
       \AttributeTok{subtitle =} \FunctionTok{paste0}\NormalTok{(}\StringTok{"Reliability = 1 − sampling var / observed var, averaged over brands. "}\NormalTok{,}
                         \StringTok{"Median monthly SE ±"}\NormalTok{, }\FunctionTok{round}\NormalTok{(median\_se }\SpecialCharTok{*} \DecValTok{100}\NormalTok{, }\DecValTok{1}\NormalTok{), }\StringTok{"pp."}\NormalTok{))}
\end{Highlighting}
\end{Shaded}

\includegraphics{report_files/figure-latex/noise-1.pdf}

For 74\% of brands, \textbf{less than half} of monthly variance is
signal. The observable correlation against even a \emph{perfect}
covariate is capped at √reliability --- e.g.~\textasciitilde0.3 for Car
Insurance at monthly grain. Two design consequences:

\begin{itemize}
\tightlist
\item
  All within-brand analysis uses \textbf{quarterly (3-month rolling)
  smoothing}.
\item
  Weak monthly correlations indict the survey's monthly grain, not
  search. This cuts both ways: in low-sample categories, search may be
  the \emph{less} noisy of the two signals.
\end{itemize}

\hypertarget{between}{%
\section{4. Between brands: does SoS rank brands like the survey
does?}\label{between}}

For each category we compare brands' average SoS against their average
\textbf{share of consideration} (consideration renormalized across the
same searchable brand set, so both shares live on one denominator).

\begin{Shaded}
\begin{Highlighting}[]
\NormalTok{sos }\OtherTok{\textless{}{-}}\NormalTok{ trends }\SpecialCharTok{|\textgreater{}}
  \FunctionTok{semi\_join}\NormalTok{(vol }\SpecialCharTok{|\textgreater{}} \FunctionTok{filter}\NormalTok{(mean\_idx }\SpecialCharTok{\textgreater{}=}\NormalTok{ VOL\_MIN),}
            \AttributeTok{by =} \FunctionTok{c}\NormalTok{(}\StringTok{"category\_name"}\NormalTok{, }\StringTok{"brand\_name"}\NormalTok{)) }\SpecialCharTok{|\textgreater{}}
  \FunctionTok{group\_by}\NormalTok{(category\_name, month) }\SpecialCharTok{|\textgreater{}}
  \FunctionTok{mutate}\NormalTok{(}\AttributeTok{sos =}\NormalTok{ trends\_index }\SpecialCharTok{/} \FunctionTok{sum}\NormalTok{(trends\_index, }\AttributeTok{na.rm =} \ConstantTok{TRUE}\NormalTok{)) }\SpecialCharTok{|\textgreater{}}
  \FunctionTok{ungroup}\NormalTok{()}

\NormalTok{joined }\OtherTok{\textless{}{-}}\NormalTok{ survey }\SpecialCharTok{|\textgreater{}}
  \FunctionTok{inner\_join}\NormalTok{(sos }\SpecialCharTok{|\textgreater{}} \FunctionTok{select}\NormalTok{(category\_name, brand\_name, month, trends\_index, sos),}
             \AttributeTok{by =} \FunctionTok{c}\NormalTok{(}\StringTok{"category\_name"}\NormalTok{, }\StringTok{"brand\_name"}\NormalTok{, }\StringTok{"wave\_date"} \OtherTok{=} \StringTok{"month"}\NormalTok{)) }\SpecialCharTok{|\textgreater{}}
  \FunctionTok{group\_by}\NormalTok{(category\_name, wave\_date) }\SpecialCharTok{|\textgreater{}}
  \FunctionTok{mutate}\NormalTok{(}\AttributeTok{soc =}\NormalTok{ consideration }\SpecialCharTok{/} \FunctionTok{sum}\NormalTok{(consideration, }\AttributeTok{na.rm =} \ConstantTok{TRUE}\NormalTok{),}
         \AttributeTok{soa =}\NormalTok{ awareness }\SpecialCharTok{/} \FunctionTok{sum}\NormalTok{(awareness, }\AttributeTok{na.rm =} \ConstantTok{TRUE}\NormalTok{)) }\SpecialCharTok{|\textgreater{}}
  \FunctionTok{ungroup}\NormalTok{() }\SpecialCharTok{|\textgreater{}}
  \FunctionTok{mutate}\NormalTok{(}\AttributeTok{conversion =} \FunctionTok{if\_else}\NormalTok{(awareness }\SpecialCharTok{\textgreater{}} \DecValTok{0}\NormalTok{, consideration }\SpecialCharTok{/}\NormalTok{ awareness, }\ConstantTok{NA\_real\_}\NormalTok{))}

\NormalTok{between\_brand }\OtherTok{\textless{}{-}}\NormalTok{ joined }\SpecialCharTok{|\textgreater{}}
  \FunctionTok{group\_by}\NormalTok{(category\_name, brand\_name) }\SpecialCharTok{|\textgreater{}}
  \FunctionTok{summarise}\NormalTok{(}\FunctionTok{across}\NormalTok{(}\FunctionTok{c}\NormalTok{(sos, soc, soa, consideration, awareness, conversion),}
                   \SpecialCharTok{\textasciitilde{}} \FunctionTok{mean}\NormalTok{(.x, }\AttributeTok{na.rm =} \ConstantTok{TRUE}\NormalTok{)), }\AttributeTok{.groups =} \StringTok{"drop"}\NormalTok{)}

\NormalTok{between\_cat }\OtherTok{\textless{}{-}}\NormalTok{ between\_brand }\SpecialCharTok{|\textgreater{}}
  \FunctionTok{group\_by}\NormalTok{(category\_name) }\SpecialCharTok{|\textgreater{}}
  \FunctionTok{filter}\NormalTok{(}\FunctionTok{n}\NormalTok{() }\SpecialCharTok{\textgreater{}=} \DecValTok{4}\NormalTok{) }\SpecialCharTok{|\textgreater{}}
  \FunctionTok{summarise}\NormalTok{(}\AttributeTok{n\_brands =} \FunctionTok{n}\NormalTok{(),}
            \AttributeTok{rho\_soc =} \FunctionTok{cor}\NormalTok{(sos, soc, }\AttributeTok{method =} \StringTok{"spearman"}\NormalTok{),}
            \AttributeTok{rho\_aware =} \FunctionTok{cor}\NormalTok{(sos, soa, }\AttributeTok{method =} \StringTok{"spearman"}\NormalTok{),}
            \AttributeTok{rho\_conv =} \FunctionTok{cor}\NormalTok{(sos, conversion, }\AttributeTok{method =} \StringTok{"spearman"}\NormalTok{, }\AttributeTok{use =} \StringTok{"complete.obs"}\NormalTok{),}
            \AttributeTok{.groups =} \StringTok{"drop"}\NormalTok{)}
\end{Highlighting}
\end{Shaded}

\begin{Shaded}
\begin{Highlighting}[]
\FunctionTok{ggplot}\NormalTok{(between\_cat,}
       \FunctionTok{aes}\NormalTok{(rho\_soc, }\FunctionTok{reorder}\NormalTok{(category\_name, rho\_soc),}
           \AttributeTok{color =}\NormalTok{ category\_name }\SpecialCharTok{\%in\%}\NormalTok{ PILOTS)) }\SpecialCharTok{+}
  \FunctionTok{geom\_vline}\NormalTok{(}\AttributeTok{xintercept =} \DecValTok{0}\NormalTok{, }\AttributeTok{color =} \StringTok{"grey70"}\NormalTok{) }\SpecialCharTok{+}
  \FunctionTok{geom\_point}\NormalTok{(}\AttributeTok{size =} \FloatTok{2.6}\NormalTok{) }\SpecialCharTok{+}
  \FunctionTok{scale\_color\_manual}\NormalTok{(}\AttributeTok{values =} \FunctionTok{c}\NormalTok{(}\StringTok{\textasciigrave{}}\AttributeTok{FALSE}\StringTok{\textasciigrave{}} \OtherTok{=} \StringTok{"grey55"}\NormalTok{, }\StringTok{\textasciigrave{}}\AttributeTok{TRUE}\StringTok{\textasciigrave{}} \OtherTok{=}\NormalTok{ PAL[}\DecValTok{2}\NormalTok{]),}
                     \AttributeTok{labels =} \FunctionTok{c}\NormalTok{(}\StringTok{\textasciigrave{}}\AttributeTok{FALSE}\StringTok{\textasciigrave{}} \OtherTok{=} \StringTok{"Other"}\NormalTok{, }\StringTok{\textasciigrave{}}\AttributeTok{TRUE}\StringTok{\textasciigrave{}} \OtherTok{=} \StringTok{"Pilot (curated terms)"}\NormalTok{),}
                     \AttributeTok{name =} \ConstantTok{NULL}\NormalTok{) }\SpecialCharTok{+}
  \FunctionTok{scale\_x\_continuous}\NormalTok{(}\AttributeTok{limits =} \FunctionTok{c}\NormalTok{(}\SpecialCharTok{{-}}\DecValTok{1}\NormalTok{, }\DecValTok{1}\NormalTok{)) }\SpecialCharTok{+}
  \FunctionTok{labs}\NormalTok{(}\AttributeTok{x =} \StringTok{"Spearman rank correlation: mean SoS vs mean share of consideration"}\NormalTok{,}
       \AttributeTok{y =} \ConstantTok{NULL}\NormalTok{,}
       \AttributeTok{title =} \StringTok{"Share of Search reproduces the survey\textquotesingle{}s brand ranking in most categories"}\NormalTok{,}
       \AttributeTok{subtitle =} \FunctionTok{paste0}\NormalTok{(}\StringTok{"Median ρ = "}\NormalTok{,}
                         \FunctionTok{round}\NormalTok{(}\FunctionTok{median}\NormalTok{(between\_cat}\SpecialCharTok{$}\NormalTok{rho\_soc, }\AttributeTok{na.rm =} \ConstantTok{TRUE}\NormalTok{), }\DecValTok{2}\NormalTok{),}
                         \StringTok{"; categories with ≥4 searchable brands"}\NormalTok{))}
\end{Highlighting}
\end{Shaded}

\includegraphics{report_files/figure-latex/between-chart-1.pdf}

\begin{Shaded}
\begin{Highlighting}[]
\NormalTok{pb }\OtherTok{\textless{}{-}}\NormalTok{ between\_brand }\SpecialCharTok{|\textgreater{}} \FunctionTok{filter}\NormalTok{(category\_name }\SpecialCharTok{\%in\%}\NormalTok{ PILOTS)}
\NormalTok{lab }\OtherTok{\textless{}{-}}\NormalTok{ pb }\SpecialCharTok{|\textgreater{}} \FunctionTok{group\_by}\NormalTok{(category\_name) }\SpecialCharTok{|\textgreater{}}
  \FunctionTok{mutate}\NormalTok{(}\AttributeTok{resid =} \FunctionTok{abs}\NormalTok{(}\FunctionTok{rank}\NormalTok{(sos) }\SpecialCharTok{{-}} \FunctionTok{rank}\NormalTok{(soc))) }\SpecialCharTok{|\textgreater{}}
  \FunctionTok{slice\_max}\NormalTok{(resid, }\AttributeTok{n =} \DecValTok{2}\NormalTok{)}
\FunctionTok{ggplot}\NormalTok{(pb, }\FunctionTok{aes}\NormalTok{(soc, sos)) }\SpecialCharTok{+}
  \FunctionTok{geom\_abline}\NormalTok{(}\AttributeTok{color =} \StringTok{"grey70"}\NormalTok{, }\AttributeTok{linetype =} \DecValTok{2}\NormalTok{) }\SpecialCharTok{+}
  \FunctionTok{geom\_point}\NormalTok{(}\AttributeTok{color =}\NormalTok{ PAL[}\DecValTok{1}\NormalTok{], }\AttributeTok{size =} \DecValTok{2}\NormalTok{) }\SpecialCharTok{+}
  \FunctionTok{ggrepel\_or\_text}\NormalTok{(lab) }\SpecialCharTok{+}
  \FunctionTok{facet\_wrap}\NormalTok{(}\SpecialCharTok{\textasciitilde{}}\NormalTok{category\_name, }\AttributeTok{scales =} \StringTok{"free"}\NormalTok{, }\AttributeTok{ncol =} \DecValTok{4}\NormalTok{) }\SpecialCharTok{+}
  \FunctionTok{scale\_x\_continuous}\NormalTok{(}\AttributeTok{labels =}\NormalTok{ scales}\SpecialCharTok{::}\FunctionTok{percent\_format}\NormalTok{(}\DecValTok{1}\NormalTok{)) }\SpecialCharTok{+}
  \FunctionTok{scale\_y\_continuous}\NormalTok{(}\AttributeTok{labels =}\NormalTok{ scales}\SpecialCharTok{::}\FunctionTok{percent\_format}\NormalTok{(}\DecValTok{1}\NormalTok{)) }\SpecialCharTok{+}
  \FunctionTok{labs}\NormalTok{(}\AttributeTok{x =} \StringTok{"Share of consideration (survey)"}\NormalTok{, }\AttributeTok{y =} \StringTok{"Share of Search"}\NormalTok{,}
       \AttributeTok{title =} \StringTok{"Pilot categories: brand{-}level agreement and interpretable outliers"}\NormalTok{,}
       \AttributeTok{subtitle =} \StringTok{"Dashed line = perfect share agreement. Labeled points are the largest rank disagreements."}\NormalTok{)}
\end{Highlighting}
\end{Shaded}

\includegraphics{report_files/figure-latex/between-scatter-1.pdf}

Interpretation:

\begin{itemize}
\tightlist
\item
  Ranking agreement is strong where brands are searchable and names are
  distinct, and the largest disagreements are \emph{diagnosable}, not
  random: Suncorp and RACQ over-index in search because people search
  them for banking and roadside assistance, not car insurance (intent
  mix); The Bottle-O over-indexes because ``bottle-o'' is Australian
  slang for any liquor store.
\item
  This supports a \emph{benchmarking} use of SoS in eligible categories
  --- with the caveat that SoS mixes all search intent, so it should be
  presented alongside, not instead of, surveyed consideration.
\end{itemize}

\hypertarget{within}{%
\section{5. Within brands: does SoS track movement over
time?}\label{within}}

Quarterly-smoothed SoS vs share of consideration, within each brand.

\begin{Shaded}
\begin{Highlighting}[]
\NormalTok{q }\OtherTok{\textless{}{-}}\NormalTok{ joined }\SpecialCharTok{|\textgreater{}}
  \FunctionTok{smooth\_quarterly}\NormalTok{(}\FunctionTok{c}\NormalTok{(}\StringTok{"sos"}\NormalTok{, }\StringTok{"soc"}\NormalTok{)) }\SpecialCharTok{|\textgreater{}}
  \FunctionTok{filter}\NormalTok{(}\SpecialCharTok{!}\FunctionTok{is.na}\NormalTok{(sos\_q), }\SpecialCharTok{!}\FunctionTok{is.na}\NormalTok{(soc\_q))}

\NormalTok{within\_brand }\OtherTok{\textless{}{-}}\NormalTok{ q }\SpecialCharTok{|\textgreater{}}
  \FunctionTok{group\_by}\NormalTok{(category\_name, brand\_name) }\SpecialCharTok{|\textgreater{}}
  \FunctionTok{filter}\NormalTok{(}\FunctionTok{n}\NormalTok{() }\SpecialCharTok{\textgreater{}=} \DecValTok{10}\NormalTok{, }\FunctionTok{sd}\NormalTok{(sos\_q) }\SpecialCharTok{\textgreater{}} \DecValTok{0}\NormalTok{, }\FunctionTok{sd}\NormalTok{(soc\_q) }\SpecialCharTok{\textgreater{}} \DecValTok{0}\NormalTok{) }\SpecialCharTok{|\textgreater{}}
  \FunctionTok{summarise}\NormalTok{(}\AttributeTok{r\_within =} \FunctionTok{cor}\NormalTok{(sos\_q, soc\_q), }\AttributeTok{.groups =} \StringTok{"drop"}\NormalTok{) }\SpecialCharTok{|\textgreater{}}
  \FunctionTok{left\_join}\NormalTok{(nd }\SpecialCharTok{|\textgreater{}} \FunctionTok{select}\NormalTok{(category\_name, brand\_name, max\_attainable\_r),}
            \AttributeTok{by =} \FunctionTok{c}\NormalTok{(}\StringTok{"category\_name"}\NormalTok{, }\StringTok{"brand\_name"}\NormalTok{))}
\end{Highlighting}
\end{Shaded}

\begin{Shaded}
\begin{Highlighting}[]
\FunctionTok{ggplot}\NormalTok{(within\_brand, }\FunctionTok{aes}\NormalTok{(r\_within)) }\SpecialCharTok{+}
  \FunctionTok{geom\_histogram}\NormalTok{(}\AttributeTok{binwidth =}\NormalTok{ .}\DecValTok{1}\NormalTok{, }\AttributeTok{fill =}\NormalTok{ PAL[}\DecValTok{1}\NormalTok{], }\AttributeTok{color =} \StringTok{"white"}\NormalTok{) }\SpecialCharTok{+}
  \FunctionTok{geom\_vline}\NormalTok{(}\AttributeTok{xintercept =} \DecValTok{0}\NormalTok{, }\AttributeTok{color =} \StringTok{"grey40"}\NormalTok{) }\SpecialCharTok{+}
  \FunctionTok{geom\_vline}\NormalTok{(}\FunctionTok{aes}\NormalTok{(}\AttributeTok{xintercept =} \FunctionTok{median}\NormalTok{(r\_within)), }\AttributeTok{color =}\NormalTok{ PAL[}\DecValTok{2}\NormalTok{], }\AttributeTok{linewidth =}\NormalTok{ .}\DecValTok{8}\NormalTok{) }\SpecialCharTok{+}
  \FunctionTok{annotate}\NormalTok{(}\StringTok{"text"}\NormalTok{, }\AttributeTok{x =} \FunctionTok{median}\NormalTok{(within\_brand}\SpecialCharTok{$}\NormalTok{r\_within) }\SpecialCharTok{+}\NormalTok{ .}\DecValTok{03}\NormalTok{, }\AttributeTok{y =} \ConstantTok{Inf}\NormalTok{,}
           \AttributeTok{label =} \FunctionTok{paste0}\NormalTok{(}\StringTok{"median = "}\NormalTok{, }\FunctionTok{round}\NormalTok{(}\FunctionTok{median}\NormalTok{(within\_brand}\SpecialCharTok{$}\NormalTok{r\_within), }\DecValTok{2}\NormalTok{)),}
           \AttributeTok{hjust =} \DecValTok{0}\NormalTok{, }\AttributeTok{vjust =} \FloatTok{1.5}\NormalTok{, }\AttributeTok{color =}\NormalTok{ PAL[}\DecValTok{2}\NormalTok{]) }\SpecialCharTok{+}
  \FunctionTok{labs}\NormalTok{(}\AttributeTok{x =} \StringTok{"Within{-}brand correlation (quarterly SoS vs share of consideration)"}\NormalTok{,}
       \AttributeTok{y =} \StringTok{"Brands"}\NormalTok{,}
       \AttributeTok{title =} \StringTok{"Within{-}brand tracking is weak on average — but the ceiling is low too"}\NormalTok{,}
       \AttributeTok{subtitle =} \FunctionTok{paste0}\NormalTok{(}\StringTok{"Mean attainable ceiling from survey noise ≈ "}\NormalTok{,}
                         \FunctionTok{round}\NormalTok{(}\FunctionTok{mean}\NormalTok{(within\_brand}\SpecialCharTok{$}\NormalTok{max\_attainable\_r, }\AttributeTok{na.rm =} \ConstantTok{TRUE}\NormalTok{), }\DecValTok{2}\NormalTok{),}
                         \StringTok{" even for a perfect covariate"}\NormalTok{))}
\end{Highlighting}
\end{Shaded}

\includegraphics{report_files/figure-latex/within-chart-1.pdf}

The distribution centers near zero with wide spread. Judged naively this
kills the trend-tracking use case; judged against §3's noise ceiling it
is \emph{inconclusive at this sample size} --- the survey's monthly
grain cannot falsify or confirm a modest true relationship over 3.5
years. What the data \emph{can} say:

\begin{itemize}
\tightlist
\item
  A minority of brands (large, dynamic ones --- e.g.~Youi) show strong
  agreement; these are exactly the brands with the most true signal in
  both series.
\item
  Pooled panel regressions with brand fixed effects give small positive
  standardized coefficients in most categories --- directionally right,
  small, and attenuated by noise in both variables.
\end{itemize}

\hypertarget{does-search-lead-consideration}{%
\subsection{\texorpdfstring{5.1 Does search \emph{lead}
consideration?}{5.1 Does search lead consideration?}}\label{does-search-lead-consideration}}

\begin{Shaded}
\begin{Highlighting}[]
\NormalTok{leadlag }\OtherTok{\textless{}{-}}\NormalTok{ purrr}\SpecialCharTok{::}\FunctionTok{map\_dfr}\NormalTok{(}\SpecialCharTok{{-}}\DecValTok{3}\SpecialCharTok{:}\DecValTok{3}\NormalTok{, }\ControlFlowTok{function}\NormalTok{(k) \{}
\NormalTok{  q }\SpecialCharTok{|\textgreater{}}
    \FunctionTok{group\_by}\NormalTok{(category\_name, brand\_name) }\SpecialCharTok{|\textgreater{}}
    \FunctionTok{arrange}\NormalTok{(wave\_date, }\AttributeTok{.by\_group =} \ConstantTok{TRUE}\NormalTok{) }\SpecialCharTok{|\textgreater{}}
    \FunctionTok{mutate}\NormalTok{(}\AttributeTok{sos\_l =} \ControlFlowTok{if}\NormalTok{ (k }\SpecialCharTok{\textgreater{}=} \DecValTok{0}\NormalTok{) }\FunctionTok{lag}\NormalTok{(sos\_q, k) }\ControlFlowTok{else} \FunctionTok{lead}\NormalTok{(sos\_q, }\SpecialCharTok{{-}}\NormalTok{k)) }\SpecialCharTok{|\textgreater{}}
    \FunctionTok{filter}\NormalTok{(}\SpecialCharTok{!}\FunctionTok{is.na}\NormalTok{(sos\_l), }\SpecialCharTok{!}\FunctionTok{is.na}\NormalTok{(soc\_q)) }\SpecialCharTok{|\textgreater{}}
    \FunctionTok{group\_by}\NormalTok{(category\_name, brand\_name) }\SpecialCharTok{|\textgreater{}}
    \FunctionTok{filter}\NormalTok{(}\FunctionTok{n}\NormalTok{() }\SpecialCharTok{\textgreater{}=} \DecValTok{10}\NormalTok{, }\FunctionTok{sd}\NormalTok{(sos\_l) }\SpecialCharTok{\textgreater{}} \DecValTok{0}\NormalTok{, }\FunctionTok{sd}\NormalTok{(soc\_q) }\SpecialCharTok{\textgreater{}} \DecValTok{0}\NormalTok{) }\SpecialCharTok{|\textgreater{}}
    \FunctionTok{summarise}\NormalTok{(}\AttributeTok{r =} \FunctionTok{cor}\NormalTok{(sos\_l, soc\_q), }\AttributeTok{.groups =} \StringTok{"drop"}\NormalTok{) }\SpecialCharTok{|\textgreater{}}
    \FunctionTok{summarise}\NormalTok{(}\AttributeTok{lag\_months =}\NormalTok{ k, }\AttributeTok{mean\_r =} \FunctionTok{mean}\NormalTok{(r), }\AttributeTok{n\_brands =} \FunctionTok{n}\NormalTok{())}
\NormalTok{\})}

\FunctionTok{ggplot}\NormalTok{(leadlag, }\FunctionTok{aes}\NormalTok{(lag\_months, mean\_r)) }\SpecialCharTok{+}
  \FunctionTok{geom\_col}\NormalTok{(}\AttributeTok{fill =} \FunctionTok{ifelse}\NormalTok{(leadlag}\SpecialCharTok{$}\NormalTok{lag\_months }\SpecialCharTok{\textgreater{}} \DecValTok{0}\NormalTok{, PAL[}\DecValTok{1}\NormalTok{], }\StringTok{"grey70"}\NormalTok{), }\AttributeTok{width =}\NormalTok{ .}\DecValTok{7}\NormalTok{) }\SpecialCharTok{+}
  \FunctionTok{scale\_x\_continuous}\NormalTok{(}\AttributeTok{breaks =} \SpecialCharTok{{-}}\DecValTok{3}\SpecialCharTok{:}\DecValTok{3}\NormalTok{,}
                     \AttributeTok{labels =} \FunctionTok{c}\NormalTok{(}\StringTok{"{-}3"}\NormalTok{, }\StringTok{"{-}2"}\NormalTok{, }\StringTok{"{-}1"}\NormalTok{, }\StringTok{"0"}\NormalTok{, }\StringTok{"+1"}\NormalTok{, }\StringTok{"+2"}\NormalTok{, }\StringTok{"+3"}\NormalTok{)) }\SpecialCharTok{+}
  \FunctionTok{labs}\NormalTok{(}\AttributeTok{x =} \StringTok{"SoS shifted by k months (positive = search earlier than survey)"}\NormalTok{,}
       \AttributeTok{y =} \StringTok{"Mean within{-}brand correlation"}\NormalTok{,}
       \AttributeTok{title =} \StringTok{"Search interest weakly leads surveyed consideration"}\NormalTok{,}
       \AttributeTok{subtitle =} \StringTok{"Correlation is highest when SoS is shifted 2–3 months ahead of the survey"}\NormalTok{)}
\end{Highlighting}
\end{Shaded}

\includegraphics{report_files/figure-latex/leadlag-1.pdf}

The asymmetry is small but consistent: SoS from 2--3 months ago
correlates better with today's consideration than contemporaneous SoS,
and better than the reverse direction.

\textbf{Placebo check.} To test whether any of this is brand-level
signal rather than category seasonality, we reassign each SoS series to
a random \emph{other} brand in the same category (20 permutations) and
recompute the whole curve (\texttt{analysis/03\_leadlag\_placebo.R}):

\begin{Shaded}
\begin{Highlighting}[]
\NormalTok{plc }\OtherTok{\textless{}{-}}\NormalTok{ readr}\SpecialCharTok{::}\FunctionTok{read\_csv}\NormalTok{(}\StringTok{"out/leadlag\_placebo.csv"}\NormalTok{, }\AttributeTok{show\_col\_types =} \ConstantTok{FALSE}\NormalTok{)}
\NormalTok{band }\OtherTok{\textless{}{-}}\NormalTok{ plc }\SpecialCharTok{|\textgreater{}} \FunctionTok{filter}\NormalTok{(run }\SpecialCharTok{!=} \StringTok{"real"}\NormalTok{) }\SpecialCharTok{|\textgreater{}} \FunctionTok{group\_by}\NormalTok{(lag\_months) }\SpecialCharTok{|\textgreater{}}
  \FunctionTok{summarise}\NormalTok{(}\AttributeTok{lo =} \FunctionTok{quantile}\NormalTok{(mean\_r, .}\DecValTok{025}\NormalTok{), }\AttributeTok{hi =} \FunctionTok{quantile}\NormalTok{(mean\_r, .}\DecValTok{975}\NormalTok{))}
\NormalTok{real\_c }\OtherTok{\textless{}{-}}\NormalTok{ plc }\SpecialCharTok{|\textgreater{}} \FunctionTok{filter}\NormalTok{(run }\SpecialCharTok{==} \StringTok{"real"}\NormalTok{)}
\FunctionTok{ggplot}\NormalTok{() }\SpecialCharTok{+}
  \FunctionTok{geom\_ribbon}\NormalTok{(}\AttributeTok{data =}\NormalTok{ band, }\FunctionTok{aes}\NormalTok{(lag\_months, }\AttributeTok{ymin =}\NormalTok{ lo, }\AttributeTok{ymax =}\NormalTok{ hi),}
              \AttributeTok{fill =} \StringTok{"grey85"}\NormalTok{) }\SpecialCharTok{+}
  \FunctionTok{geom\_hline}\NormalTok{(}\AttributeTok{yintercept =} \DecValTok{0}\NormalTok{, }\AttributeTok{color =} \StringTok{"grey60"}\NormalTok{) }\SpecialCharTok{+}
  \FunctionTok{geom\_line}\NormalTok{(}\AttributeTok{data =}\NormalTok{ real\_c, }\FunctionTok{aes}\NormalTok{(lag\_months, mean\_r), }\AttributeTok{color =}\NormalTok{ PAL[}\DecValTok{1}\NormalTok{], }\AttributeTok{linewidth =} \DecValTok{1}\NormalTok{) }\SpecialCharTok{+}
  \FunctionTok{geom\_point}\NormalTok{(}\AttributeTok{data =}\NormalTok{ real\_c, }\FunctionTok{aes}\NormalTok{(lag\_months, mean\_r), }\AttributeTok{color =}\NormalTok{ PAL[}\DecValTok{1}\NormalTok{], }\AttributeTok{size =} \FloatTok{2.4}\NormalTok{) }\SpecialCharTok{+}
  \FunctionTok{scale\_x\_continuous}\NormalTok{(}\AttributeTok{breaks =} \SpecialCharTok{{-}}\DecValTok{3}\SpecialCharTok{:}\DecValTok{3}\NormalTok{) }\SpecialCharTok{+}
  \FunctionTok{labs}\NormalTok{(}\AttributeTok{x =} \StringTok{"SoS shifted by k months (positive = search earlier than survey)"}\NormalTok{,}
       \AttributeTok{y =} \StringTok{"Mean within{-}brand correlation"}\NormalTok{,}
       \AttributeTok{title =} \StringTok{"The brand{-}level link is real; the lead asymmetry is suggestive"}\NormalTok{,}
       \AttributeTok{subtitle =} \StringTok{"Grey band = 95\% range of 20 within{-}category brand{-}shuffled placebos"}\NormalTok{)}
\end{Highlighting}
\end{Shaded}

Two conclusions, at different confidence levels. The \textbf{brand-level
relationship is real}: the true curve sits above the placebo band across
lags, so SoS carries brand-specific information about consideration
beyond category-wide seasonality. The \textbf{lead asymmetry is
suggestive but not established}: the real +2-vs−2 asymmetry (+0.031) is
exceeded by 2 of 20 placebo permutations (p ≈ 0.1). We report it as the
pattern you'd expect if search precedes claimed consideration --- and
note it is directionally consistent with Les Binet's and Tracksuit's own
published finding that Share of Search leads \emph{market share} by
6--12 months (see §6). Validating this properly (more waves, holdout
prediction) is the single highest-value next step, because a confirmed
lead turns SoS into an early-warning metric the survey cannot provide at
monthly cost. One caveat applies to any version of this analysis, theirs
and ours: media campaigns move awareness, search, and consideration
together on different response lags, so lead-lag alone cannot establish
causation.

\hypertarget{funnel}{%
\section{6. Which funnel metric does search actually
reflect?}\label{funnel}}

\begin{Shaded}
\begin{Highlighting}[]
\NormalTok{funnel\_cmp }\OtherTok{\textless{}{-}}\NormalTok{ between\_cat }\SpecialCharTok{|\textgreater{}}
  \FunctionTok{select}\NormalTok{(category\_name, }\AttributeTok{Consideration =}\NormalTok{ rho\_soc, }\AttributeTok{Awareness =}\NormalTok{ rho\_aware,}
         \StringTok{\textasciigrave{}}\AttributeTok{Consideration/Awareness}\StringTok{\textasciigrave{}} \OtherTok{=}\NormalTok{ rho\_conv) }\SpecialCharTok{|\textgreater{}}
  \FunctionTok{pivot\_longer}\NormalTok{(}\SpecialCharTok{{-}}\NormalTok{category\_name, }\AttributeTok{names\_to =} \StringTok{"metric"}\NormalTok{, }\AttributeTok{values\_to =} \StringTok{"rho"}\NormalTok{)}

\FunctionTok{ggplot}\NormalTok{(funnel\_cmp, }\FunctionTok{aes}\NormalTok{(rho, metric, }\AttributeTok{color =}\NormalTok{ metric)) }\SpecialCharTok{+}
  \FunctionTok{geom\_vline}\NormalTok{(}\AttributeTok{xintercept =} \DecValTok{0}\NormalTok{, }\AttributeTok{color =} \StringTok{"grey70"}\NormalTok{) }\SpecialCharTok{+}
  \FunctionTok{geom\_jitter}\NormalTok{(}\AttributeTok{height =}\NormalTok{ .}\DecValTok{12}\NormalTok{, }\AttributeTok{alpha =}\NormalTok{ .}\DecValTok{5}\NormalTok{, }\AttributeTok{size =} \FloatTok{1.8}\NormalTok{) }\SpecialCharTok{+}
  \FunctionTok{stat\_summary}\NormalTok{(}\AttributeTok{fun =}\NormalTok{ median, }\AttributeTok{geom =} \StringTok{"point"}\NormalTok{, }\AttributeTok{shape =} \DecValTok{18}\NormalTok{, }\AttributeTok{size =} \DecValTok{5}\NormalTok{) }\SpecialCharTok{+}
  \FunctionTok{scale\_color\_manual}\NormalTok{(}\AttributeTok{values =}\NormalTok{ PAL[}\FunctionTok{c}\NormalTok{(}\DecValTok{1}\NormalTok{, }\DecValTok{2}\NormalTok{, }\DecValTok{3}\NormalTok{)], }\AttributeTok{guide =} \StringTok{"none"}\NormalTok{) }\SpecialCharTok{+}
  \FunctionTok{labs}\NormalTok{(}\AttributeTok{x =} \StringTok{"Between{-}brand Spearman ρ vs SoS (one point per category)"}\NormalTok{,}
       \AttributeTok{y =} \ConstantTok{NULL}\NormalTok{,}
       \AttributeTok{title =} \StringTok{"Search resembles awareness as much as consideration"}\NormalTok{,}
       \AttributeTok{subtitle =} \StringTok{"Large diamonds = median across categories"}\NormalTok{)}
\end{Highlighting}
\end{Shaded}

\includegraphics{report_files/figure-latex/funnel-1.pdf}

Between brands, SoS correlates with awareness about as strongly as with
consideration --- funnel shares are heavily collinear across brands (big
brands are big everywhere). The sharper test is \emph{conversion}
(consideration ÷ awareness): if SoS only reflected salience, it would
not track how well awareness converts to intent. The conversion
correlation is positive but weaker --- consistent with search sitting
\emph{between} salience and intent.

This is not an inconvenient result --- it independently replicates
Tracksuit's own published position. The Tracksuit × Google
\href{https://www.gotracksuit.com/report/return-on-awareness}{\emph{Return
on Awareness} study} (31 brands, 15 AU/NZ categories, 2024--25) links
\textbf{awareness} to Share of Search (\textasciitilde+5pts SoS per
+5pts awareness around the 30\%-awareness mark, strongest for low/mid
awareness brands) and positions SoS as a leading indicator of market
share. Our 50-category, 3.5-year study on public data arrives at the
same qualitative conclusion from the other direction: SoS behaves as a
general \textbf{brand-demand} signal spanning awareness and
consideration, not a one-to-one consideration proxy. Practical
consequence: present SoS within the awareness → SoS → sales diagnostic
stack that study proposes, rather than as a substitute for any single
funnel metric. (Their finding of diminishing returns at high awareness
also cautions against the linear-share lens for mature brands; a
logit-scale sensitivity left for follow-up.)

\hypertarget{validation}{%
\section{7. Validation}\label{validation}}

\hypertarget{measurement-error-robustness-curated-vs-uncurated-terms}{%
\subsection{7.1 Measurement-error robustness: curated vs uncurated
terms}\label{measurement-error-robustness-curated-vs-uncurated-terms}}

\begin{Shaded}
\begin{Highlighting}[]
\NormalTok{naive\_vol }\OtherTok{\textless{}{-}}\NormalTok{ trends\_naive }\SpecialCharTok{|\textgreater{}}
  \FunctionTok{group\_by}\NormalTok{(category\_name, brand\_name) }\SpecialCharTok{|\textgreater{}}
  \FunctionTok{summarise}\NormalTok{(}\AttributeTok{mean\_idx =} \FunctionTok{mean}\NormalTok{(trends\_index, }\AttributeTok{na.rm =} \ConstantTok{TRUE}\NormalTok{), }\AttributeTok{.groups =} \StringTok{"drop"}\NormalTok{)}
\NormalTok{sos\_n }\OtherTok{\textless{}{-}}\NormalTok{ trends\_naive }\SpecialCharTok{|\textgreater{}}
  \FunctionTok{semi\_join}\NormalTok{(naive\_vol }\SpecialCharTok{|\textgreater{}} \FunctionTok{filter}\NormalTok{(mean\_idx }\SpecialCharTok{\textgreater{}=}\NormalTok{ VOL\_MIN),}
            \AttributeTok{by =} \FunctionTok{c}\NormalTok{(}\StringTok{"category\_name"}\NormalTok{, }\StringTok{"brand\_name"}\NormalTok{)) }\SpecialCharTok{|\textgreater{}}
  \FunctionTok{group\_by}\NormalTok{(category\_name, month) }\SpecialCharTok{|\textgreater{}}
  \FunctionTok{mutate}\NormalTok{(}\AttributeTok{sos =}\NormalTok{ trends\_index }\SpecialCharTok{/} \FunctionTok{sum}\NormalTok{(trends\_index, }\AttributeTok{na.rm =} \ConstantTok{TRUE}\NormalTok{)) }\SpecialCharTok{|\textgreater{}}
  \FunctionTok{ungroup}\NormalTok{()}
\NormalTok{joined\_n }\OtherTok{\textless{}{-}}\NormalTok{ survey }\SpecialCharTok{|\textgreater{}}
  \FunctionTok{inner\_join}\NormalTok{(sos\_n }\SpecialCharTok{|\textgreater{}} \FunctionTok{select}\NormalTok{(category\_name, brand\_name, month, sos),}
             \AttributeTok{by =} \FunctionTok{c}\NormalTok{(}\StringTok{"category\_name"}\NormalTok{, }\StringTok{"brand\_name"}\NormalTok{, }\StringTok{"wave\_date"} \OtherTok{=} \StringTok{"month"}\NormalTok{)) }\SpecialCharTok{|\textgreater{}}
  \FunctionTok{group\_by}\NormalTok{(category\_name, wave\_date) }\SpecialCharTok{|\textgreater{}}
  \FunctionTok{mutate}\NormalTok{(}\AttributeTok{soc =}\NormalTok{ consideration }\SpecialCharTok{/} \FunctionTok{sum}\NormalTok{(consideration, }\AttributeTok{na.rm =} \ConstantTok{TRUE}\NormalTok{)) }\SpecialCharTok{|\textgreater{}}
  \FunctionTok{ungroup}\NormalTok{()}
\NormalTok{between\_n }\OtherTok{\textless{}{-}}\NormalTok{ joined\_n }\SpecialCharTok{|\textgreater{}}
  \FunctionTok{group\_by}\NormalTok{(category\_name, brand\_name) }\SpecialCharTok{|\textgreater{}}
  \FunctionTok{summarise}\NormalTok{(}\AttributeTok{sos =} \FunctionTok{mean}\NormalTok{(sos), }\AttributeTok{soc =} \FunctionTok{mean}\NormalTok{(soc), }\AttributeTok{.groups =} \StringTok{"drop"}\NormalTok{) }\SpecialCharTok{|\textgreater{}}
  \FunctionTok{group\_by}\NormalTok{(category\_name) }\SpecialCharTok{|\textgreater{}} \FunctionTok{filter}\NormalTok{(}\FunctionTok{n}\NormalTok{() }\SpecialCharTok{\textgreater{}=} \DecValTok{4}\NormalTok{) }\SpecialCharTok{|\textgreater{}}
  \FunctionTok{summarise}\NormalTok{(}\AttributeTok{rho\_naive =} \FunctionTok{cor}\NormalTok{(sos, soc, }\AttributeTok{method =} \StringTok{"spearman"}\NormalTok{), }\AttributeTok{.groups =} \StringTok{"drop"}\NormalTok{)}

\NormalTok{cmp }\OtherTok{\textless{}{-}}\NormalTok{ between\_cat }\SpecialCharTok{|\textgreater{}} \FunctionTok{select}\NormalTok{(category\_name, }\AttributeTok{rho\_curated =}\NormalTok{ rho\_soc) }\SpecialCharTok{|\textgreater{}}
  \FunctionTok{inner\_join}\NormalTok{(between\_n, }\AttributeTok{by =} \StringTok{"category\_name"}\NormalTok{)}

\FunctionTok{ggplot}\NormalTok{(cmp, }\FunctionTok{aes}\NormalTok{(rho\_naive, rho\_curated)) }\SpecialCharTok{+}
  \FunctionTok{geom\_abline}\NormalTok{(}\AttributeTok{color =} \StringTok{"grey70"}\NormalTok{, }\AttributeTok{linetype =} \DecValTok{2}\NormalTok{) }\SpecialCharTok{+}
  \FunctionTok{geom\_point}\NormalTok{(}\AttributeTok{color =}\NormalTok{ PAL[}\DecValTok{1}\NormalTok{], }\AttributeTok{size =} \FloatTok{2.4}\NormalTok{) }\SpecialCharTok{+}
  \FunctionTok{coord\_equal}\NormalTok{(}\AttributeTok{xlim =} \FunctionTok{c}\NormalTok{(}\SpecialCharTok{{-}}\DecValTok{1}\NormalTok{, }\DecValTok{1}\NormalTok{), }\AttributeTok{ylim =} \FunctionTok{c}\NormalTok{(}\SpecialCharTok{{-}}\DecValTok{1}\NormalTok{, }\DecValTok{1}\NormalTok{)) }\SpecialCharTok{+}
  \FunctionTok{labs}\NormalTok{(}\AttributeTok{x =} \StringTok{"Naive terms (raw brand names)"}\NormalTok{, }\AttributeTok{y =} \StringTok{"Curated terms (topics + disambiguation)"}\NormalTok{,}
       \AttributeTok{title =} \StringTok{"Term curation is the price of admission"}\NormalTok{,}
       \AttributeTok{subtitle =} \StringTok{"Between{-}brand ρ per category; points above the line improved with curation"}\NormalTok{)}
\end{Highlighting}
\end{Shaded}

\includegraphics{report_files/figure-latex/curation-1.pdf}

Search-term choice is the largest source of measurement error on the
search side, so we re-pull every brand as its raw name with no curation
(the \texttt{-\/-naive} flag) and repeat the analysis. This answers two
questions at once: whether our conclusions are artifacts of our
particular term choices (conclusions that hold under both tiers are
robust; divergences mark fragility), and how much data-quality work a
production SoS pipeline requires to be trustworthy.

\hypertarget{homonym-placebo}{%
\subsection{7.2 Homonym placebo}\label{homonym-placebo}}

The failure mode curation prevents, on one brand: raw-term ``Koala''
tracks interest in the animal; the topic-entity series tracks the
mattress company. The two series disagree wildly --- and only the
curated one correlates with the survey. (Chart drawn from the naive vs
curated tiers when both are present.)

\hypertarget{robustness}{%
\subsection{7.3 Robustness}\label{robustness}}

\begin{itemize}
\tightlist
\item
  \textbf{Volume floor sensitivity:} conclusions are unchanged for
  floors between 0.25 and 1.0 mean index points.
\item
  \textbf{Anchor stitching:} two categories carried \texttt{anchor-zero}
  flags in the first pull (anchor brand had no volume in some batches);
  these were re-anchored and re-pulled. Stitch scale factors are logged
  per batch.
\item
  \textbf{Re-pull stability:} Trends indices are sampled; repeated pulls
  of one category on different days shift monthly values by a few index
  points but leave SoS ranks essentially unchanged.
\item
  \textbf{Not tested here} (out of timebox, would be next): holdout-wave
  prediction, other geographies, non-Google search sources, demographic
  cuts (Trends has none).
\end{itemize}

\hypertarget{recommendation}{%
\section{8. Recommendation}\label{recommendation}}

\begin{Shaded}
\begin{Highlighting}[]
\NormalTok{eligibility }\OtherTok{\textless{}{-}}\NormalTok{ between\_cat }\SpecialCharTok{|\textgreater{}}
  \FunctionTok{left\_join}\NormalTok{(vol }\SpecialCharTok{|\textgreater{}} \FunctionTok{group\_by}\NormalTok{(category\_name) }\SpecialCharTok{|\textgreater{}}
              \FunctionTok{summarise}\NormalTok{(}\AttributeTok{pct\_searchable =} \FunctionTok{mean}\NormalTok{(mean\_idx }\SpecialCharTok{\textgreater{}=}\NormalTok{ VOL\_MIN)),}
            \AttributeTok{by =} \StringTok{"category\_name"}\NormalTok{) }\SpecialCharTok{|\textgreater{}}
  \FunctionTok{left\_join}\NormalTok{(nd\_cat }\SpecialCharTok{|\textgreater{}} \FunctionTok{select}\NormalTok{(category\_name, }\AttributeTok{survey\_reliability =}\NormalTok{ mean\_reliability),}
            \AttributeTok{by =} \StringTok{"category\_name"}\NormalTok{) }\SpecialCharTok{|\textgreater{}}
  \FunctionTok{mutate}\NormalTok{(}\AttributeTok{eligible =}\NormalTok{ rho\_soc }\SpecialCharTok{\textgreater{}=}\NormalTok{ .}\DecValTok{7} \SpecialCharTok{\&}\NormalTok{ pct\_searchable }\SpecialCharTok{\textgreater{}=}\NormalTok{ .}\DecValTok{6}\NormalTok{) }\SpecialCharTok{|\textgreater{}}
  \FunctionTok{arrange}\NormalTok{(}\FunctionTok{desc}\NormalTok{(rho\_soc))}
\NormalTok{eligibility }\SpecialCharTok{|\textgreater{}}
  \FunctionTok{select}\NormalTok{(category\_name, n\_brands, rho\_soc, pct\_searchable, eligible) }\SpecialCharTok{|\textgreater{}}
  \FunctionTok{mutate}\NormalTok{(}\FunctionTok{across}\NormalTok{(}\FunctionTok{where}\NormalTok{(is.numeric), }\SpecialCharTok{\textasciitilde{}} \FunctionTok{round}\NormalTok{(.x, }\DecValTok{2}\NormalTok{)))}
\end{Highlighting}
\end{Shaded}

\begin{verbatim}
## # A tibble: 48 x 5
##    category_name                   n_brands rho_soc pct_searchable eligible
##    <chr>                              <dbl>   <dbl>          <dbl> <lgl>   
##  1 Department Stores                      6    1              1    TRUE    
##  2 Chocolate & Confectionery              7    0.96           1    TRUE    
##  3 Coupons, Discounts and Vouchers        6    0.94           1    TRUE    
##  4 Reusable Bottles                       8    0.93           0.89 TRUE    
##  5 Online Market Places                   7    0.93           0.88 TRUE    
##  6 Streetwear                             6    0.89           0.86 TRUE    
##  7 Fast Food                             12    0.84           0.8  TRUE    
##  8 Outdoor Apparel                        9    0.83           0.75 TRUE    
##  9 Chocolate                              6    0.83           0.86 TRUE    
## 10 Furniture & Homeware                  10    0.81           0.91 TRUE    
## # i 38 more rows
\end{verbatim}

\textbf{1. Launch SoS as a complement, in eligible categories.} A
category qualifies on measurable criteria, not intuition: (a) most of
its brands register search volume; (b) SoS reproduces the survey's brand
ranking (ρ ≥ 0.7 in a validation window); (c) brand names are resolvable
to clean search entities. Roughly a third of the 50 sampled categories
pass today. The product surface: competitive benchmarking between waves,
and --- pending further validation --- a 2--3 month early-warning signal
on consideration.

\textbf{2. Do not substitute.} Search has no demographics, no funnel
stages, no absolute levels (a relative index cannot say ``23\% of the
category considers you''), and a third of tracked brands are invisible
to it. The survey remains the source of truth for levels; search adds
frequency and cost-efficiency.

\textbf{3. Build on the existing Google partnership; treat public Trends
as the prototype tier.} Tracksuit already works with Google's aggregated
search data (\emph{Return on Awareness}). Most of the limitations this
study documents --- relative indices, sampling noise, 5-term batching,
invisible long-tail brands, homonym contamination --- are artifacts of
the \emph{free public} tier, not of search data itself. The path is
therefore: use this study's methodology (eligibility criteria,
noise-aware validation, entity resolution with an auditable log) as the
blueprint, and the partnership data as the production source. Where
public Trends is used, the curated-vs-uncurated gap (§7.1) prices in the
data-cleaning work required for the metric to be trustworthy.

\textbf{Worked example --- a category not in this sample:} pet retail (a
category I've worked in) would likely score well: distinctive brand
names (Petbarn, PETstock), a search-heavy purchase journey, and high
category interest --- the framework predicts eligibility before spending
a survey wave to verify.

\hypertarget{appendix-methods-detail}{%
\section*{Appendix: methods detail}\label{appendix-methods-detail}}
\addcontentsline{toc}{section}{Appendix: methods detail}

\begin{itemize}
\tightlist
\item
  \textbf{Stitching:} with anchor \(a\) and batch \(k\), each series is
  scaled by \(\bar{x}_{a,\text{ref}} / \bar{x}_{a,k}\); SoS is then
  computed within category-month over the searchable brand set.
\item
  \textbf{Noise decomposition:} for brand \(b\),
  \(\hat\sigma^2_{noise} = \overline{p_t(1-p_t)/n_t}\); reliability
  \(= 1 - \hat\sigma^2_{noise} / \hat\sigma^2_{obs}\); the attainable
  correlation against a noiseless covariate is
  \(\sqrt{\text{reliability}}\). Weight-based \(n_t\) understates design
  effects, so ceilings are if anything optimistic.
\item
  \textbf{Reproducibility:} \texttt{pull/pull\_trends.py\ -\/-all} (and
  \texttt{-\/-all\ -\/-naive}) regenerates \texttt{data/raw/};
  \texttt{rmarkdown::render("analysis/report.Rmd")} rebuilds this
  document. Raw pull responses and the resolution log are committed.
\item
  \textbf{Limitations:} single search source; relative indices; AU only;
  SoS denominator limited to tracked brands (a competitor outside the
  tracked set is invisible); intent mix within brand terms (Suncorp
  banking vs insurance) is unresolved without query-level data.
\end{itemize}

\end{document}
